% Thesis update pack generated on 2026-07-08.
% Purpose: paste into Overleaf main_v2_evidence_guided.tex to update the
% software evidence, C0/C1/C2 results, Stage 8 preview, hardware deployment
% boundary, and limitations.
%
% Recommended location:
% 1. Put the C2/Stage 8 sections in Results or Evaluation.
% 2. Put the hardware deployment section near Hardware/Implementation.
% 3. Put the limitations section in Discussion/Future Work.
% 4. Import the figures listed at the bottom into Overleaf before compiling.

\section{Current Controller Evidence}
\label{sec:current_controller_evidence}

The current software pipeline evaluates three controller layers under a
matched evidence framework. C0 is a joint-reference playback baseline. C1 adds
a fixed-phase FIR-QP-MPC pitch-recovery layer. C2 keeps the C1 correction layer
and adds contact-event phase synchronisation for toe/heel liftoff and touchdown
events. This staged comparison separates baseline stability, predictive pitch
correction, and contact-timing robustness.

Table~\ref{tab:c0_c1_c2_current_evidence_20260708} summarises the current
pitch-disturbance and runtime evidence. Relative to C0, both C1 and C2 reduce
the integrated absolute pitch error in the forward and backward disturbance
cases. In the current double-support CoM-shift reference, C2 gives the same
pitch-recovery metrics as C1 because the nominal reference does not contain
scheduled contact transitions. Therefore the C2 benefit should not be claimed
from the pitch-recovery test alone; it is supported by the separate
contact-event timing tests described below.

\begin{table}[htbp]
\centering
\caption{Current C0/C1/C2 pitch-disturbance and runtime evidence.}
\label{tab:c0_c1_c2_current_evidence_20260708}
\begin{tabular}{@{}lllllll@{}}
\toprule
Controller & Direction & IAE (deg s) & Improvement & Settling (s) & Max corr. (deg) & Median solve (us) \\
\midrule
C0 & forward & 0.2710 & -- & 0.248 & 0.000 & -- \\
C1 & forward & 0.2474 & 8.7\% & 0.166 & 0.711 & 711.3 \\
C2 & forward & 0.2474 & 8.7\% & 0.166 & 0.711 & 677.2 \\
C0 & backward & 0.2681 & -- & 0.246 & 0.000 & -- \\
C1 & backward & 0.2295 & 14.4\% & 0.166 & 0.669 & 677.6 \\
C2 & backward & 0.2295 & 14.4\% & 0.166 & 0.669 & 747.2 \\
\bottomrule
\end{tabular}
\end{table}

\section{C2 Contact-Phase Synchronisation}
\label{sec:c2_contact_phase_synchronisation_20260708}

The C2 layer is implemented as a contact-event phase synchronisation mechanism
around the C1 FIR-QP-MPC controller. Its purpose is to keep the reference phase
consistent with the measured contact phase when toe/heel contact events occur
earlier or later than expected. A scripted contact-switching mechanism test
validates this event path independently from the no-transition CoM-shift pitch
test. In that mechanism test, C2 matched 2/2 scheduled contact events and the
maximum absolute phase offset remained approximately
\SI{35.0}{\milli\second}.

This result should be interpreted as contact-event mechanism validation. It
does not yet prove full FDDP walking-gait validation, because the current
walking reference is still a conservative scripted preview rather than a fully
optimised FDDP walking motion.

\section{Stage 8 Scripted Alternating-Step Preview}
\label{sec:stage8_scripted_alternating_step_20260708}

The project now includes a Stage 8 scripted alternating-step locomotion
preview. This preview is designed to exercise the left/right contact-transition
pipeline before a full FDDP walking reference is available. It contains two
scripted swing phases and scheduled left/right liftoff-touchdown events. The
test is intentionally conservative and should be reported as software
locomotion-plumbing validation rather than as final high-performance walking.

In the current MuJoCo run, Stage 8 matched 4/4 scheduled contact events,
reported a pitch RMS of \SI{0.334}{\degree}, and produced approximately
\SI{13.6}{\milli\metre} of forward displacement over
\SI{2.35}{\second}. The maximum phase offset was
\SI{48.3}{\milli\second}, and no fall was detected. The acceptance checks
confirm that all expected event kinds were observed, all expected events were
matched, no reference events were skipped, the phase-offset bound was
respected, support contact was retained, and no fall was detected.

\begin{table}[htbp]
\centering
\caption{Stage 8 scripted alternating-step preview metrics.}
\label{tab:stage8_preview_metrics_20260708}
\begin{tabular}{@{}ll@{}}
\toprule
Metric & Current value \\
\midrule
Duration & \SI{2.35}{\second} \\
Scripted swing phases & 2 \\
Scheduled contact events & 4 \\
Matched contact events & 4/4 \\
Skipped contact events & 0 \\
Maximum phase offset & \SI{48.3}{\milli\second} \\
Pitch RMS & \SI{0.334}{\degree} \\
Peak absolute pitch & \SI{0.801}{\degree} \\
Minimum torso height & \SI{0.288}{\metre} \\
Forward displacement & \SI{13.6}{\milli\metre} \\
Fall detected & False \\
\bottomrule
\end{tabular}
\end{table}

\section{Hardware Runtime Export and Deployment Boundary}
\label{sec:hardware_runtime_export_20260708}

The current hardware-side software is a MuJoCo-free Python runtime for
PC/ROS2-side deployment. It exposes C0, C1, and C2 as joint-position command
controllers for the six ST3215 joints. The exported controller package is
stored at \path{hardware/deploy/controller_package.json}. The package declares
the command type as joint positions in radians, uses a
\SI{20}{\milli\second} update period, and explicitly disallows torque
commands.

This is a controller package and runtime-loadability result, not a
microcontroller firmware image. The runtime does not directly communicate with
the ST3215 servos. A hardware bridge is still required to implement serial
communication, timeouts, emergency stop, servo fault handling, calibration, and
conversion from joint radians to calibrated servo position units.

The practical deployment architecture is therefore:

\begin{enumerate}
    \item The PC or ROS2 computer runs the C0/C1/C2 Python runtime.
    \item The runtime receives measurements such as time, pitch, pitch rate,
    and toe/heel contact states.
    \item The selected controller outputs six joint-position commands in
    radians.
    \item A hardware bridge converts these commands into ST3215-compatible
    servo commands.
    \item The control board or servo bus firmware sends the final position
    commands to the servos and reports faults or telemetry.
\end{enumerate}

The current result supports hardware readiness because the package can be
loaded and finite commands can be produced under the software smoke test. It
does not yet demonstrate real hardware walking.

\section{Simulation Coverage and Stability Conclusion}
\label{sec:simulation_coverage_stability_20260708}

The current simulation evidence covers several important but limited operating
conditions:

\begin{itemize}
    \item forward and backward CoM-shift reference replay;
    \item matched pitch-disturbance recovery for C0, C1, and C2;
    \item actuator-aware command limiting using ST3215 delay, lag, velocity,
    and rate-limit assumptions;
    \item C2 contact-switching mechanism validation;
    \item Stage 8 scripted alternating left/right contact transitions.
\end{itemize}

Within these tested conditions, the robot does not show a fall in the current
software evidence, pitch remains bounded in the Stage 8 preview, and the
controller commands remain finite and joint-limit checked in the runtime smoke
test. Therefore the current software result supports the claim that the
pipeline is stable under the tested scripted and disturbance conditions.

However, the simulations do not yet cover all conditions required to claim
general robust biped walking. A follow-on Stage 9 software robustness matrix
now covers scripted long-duration operation, friction variation,
slope-equivalent loading, small toe-step perturbations, and push pulses. In
that matrix, 14/20 scenarios pass the current acceptance checks. The stable
simulation envelope includes \SI{60}{\second} scripted operation, friction
variation from $\mu=0.50$ to $\mu=1.10$, slope-equivalent loading up to
$\pm\SI{6}{\degree}$, a \SI{2}{\milli\metre} toe step, and $\pm\SI{2}{\newton}$
push pulses. The measured failure boundary is also clear: \SI{4}{\milli\metre}
and \SI{6}{\milli\metre} toe steps fail contact-event robustness criteria,
$\pm\SI{4}{\newton}$ push pulses cause transient support contact loss, and
$\pm\SI{6}{\newton}$ push pulses cross the pitch fall threshold.

The remaining open items are therefore narrower but still important: full FDDP
walking-gait tracking, real servo communication faults, battery voltage
effects, hardware contact sensor noise, repeated hardware trial success rate,
and real hardware walking. These limitations should be stated explicitly to
avoid overclaiming.

\section{Remaining Validation Required}
\label{sec:remaining_validation_required_20260708}

The unchanged title can still be defended if the thesis clearly distinguishes
implemented evidence from next-stage validation. The remaining work is:

\begin{itemize}
    \item real hardware bring-up with calibrated ST3215 joint zero positions;
    \item a hardware bridge from C0/C1/C2 radian commands to servo bus commands;
    \item IMU pitch and pitch-rate measurement integration;
    \item toe/heel or foot-contact sensing integration;
    \item C0 hardware reference playback test;
    \item C1 hardware pitch-recovery test;
    \item C2 hardware contact-event synchronisation test;
    \item repeated trial success-rate measurement;
    \item walking duration, velocity error, step-length error, and path
    deviation metrics;
    \item battery voltage drop, servo fault events, and current/energy logs.
\end{itemize}

\section{Figures To Import}
\label{sec:figures_to_import_20260708}

The following generated figures should be imported into Overleaf and referenced
from the Results/Evaluation chapter:

\begin{itemize}
    \item \path{results/figures/c2_phase_sync_forward_pitch_comparison.png}
    \item \path{results/figures/c2_phase_sync_backward_pitch_comparison.png}
    \item \path{results/figures/c2_contact_switching_demo_phase_offset.png}
    \item \path{results/figures/c2_contact_switching_demo_contacts.png}
    \item \path{results/figures/stage8_scripted_alternating_step_contacts.png}
    \item \path{results/figures/stage8_scripted_alternating_step_phase_offset.png}
    \item \path{results/figures/stage8_scripted_alternating_step_body.png}
    \item \path{results/figures/stage9_scripted_robustness_matrix_pass_fail.png}
    \item \path{results/figures/stage9_scripted_robustness_matrix_pitch.png}
\end{itemize}
